% !TEX root = ../main.tex
\section{O pipeline de criação ponta a ponta}
\label{sec:S14-pipeline-ponta-a-ponta}

Cadeia determinística sobre artefatos em disco, ancorada no \file{.spf}.
\textbf{Regra-mãe: o renderer orquestra; o main executa.} O renderer
(\file{compilation\_flow.js}, \file{compilation\_module.js},
\file{processor\_compiler.js}) monta cada etapa como \emph{CommandSpec};
\code{runSpec} (\file{spec\_runner.js}) é o único ponto de passagem e dispara
\lstinline|exec-spec|; o main re-valida o \code{binary} contra a \emph{allowlist}
(\file{binary\_allowlist.js}) antes de spawnar. Scripts de referência no
\tool{YANC}: \file{single\_proc.sh} (um processador) e \file{multi\_proc.sh}
(vários sob top-level). Sem desvio ``tem processador?''; sem processadores os laços
são no-op (\emph{pipeline único}).

\begin{lstlisting}
=========================================================================
  PIPELINE SAPHO  (renderer orquestra | [bin] = main + allowlist)
=========================================================================
 (1) PROJETO  -> grava <proj>.spf  (fonte canonica)
     synthesizableFiles, testbenchFiles, topLevelFile,
     processors[]: name/cmmFile/clk(=100)/numClocks(=2000)/showArrays
     cria: <proj>/<proc>/{Software,Hardware,Simulation}
           components/Temp/<proc>/   (scratch -t / CWD da sim)
        v
 (2) CODIGO-FONTE  (editor Monaco)
     <proc>.cmm (ou .cpp) + top-level.v + testbench
       (<proc>_tb.v auto | top_level_tb.v | *.py cocotb)
     [Verilator] ensureChegueiToaqui -> injeta #TOAQUI no .cmm
        v
 (3) COMPILACAO  (precompileAllProcessors -> por processador)
   +-----------------------------------------------------------------+
   | [cmmcomp.exe] [-pt|-en] -i .cmm -n -p -m -t [--array]           |
   |    -> <proc>.asm + cmm_log.txt + pc_<proc>_mem.txt              |
   |       + <proc>_data.mif / <proc>_inst.mif                       |
   | [appcomp.exe] -i .asm -t  -> .asm pre-processado (params/end.)  |
   | [asmcomp.exe] -i .asm -p -d HDL -m -t -f<clk> -c<numClk>        |
   |    -> Hardware/<proc>.v (RTL) + .mif + <proc>_tb.v -> Simulation|
   +-----------------------------------------------------------------+
     (-y components/HDL: processor/core/ula/addr_dec/instr_dec/myFIFO)
        +----------------+------------------+
        v (4a Icarus)    v (4b Verilator)   v (4c cocotb/Python)
 (5) INSTRUMENTAR DUMP  (comum, ANTES de simular)
     _resolveWaveSelection:
       .gtkw ativo > Wave Config > $dumpvars no tb > default
     instrumentTestbench -> Temp/instr_<tb>.v  (copia; original intacto)
        v
   [iverilog -s tb -y HDL  [verilator --binary  [cocotb: iverilog ou
    -o .vvp]; [vvp -fst]    --timing --trace]    verilator por baixo]
     stageProcessorMemoryFiles: pc_*_mem.txt -> Temp/ ($readmemb)
        v
     DUMP de ondas (FST)
     _extractFstHeaderVcd -> [fst2vcd] -> header.vcd (escopos/sinais)
        v
 (6) VISUALIZAR  (getViewer)
     +-- gtkwave -> buildAuroraGtkw -> .gtkw -> [gtkwave.exe]
     +-- surfer  -> buildSurferLayout -> .surf -> [surfer.exe]
        faixas C+- / assembly / variaveis em LOCKSTEP com o clock
 (7) INSPECIONAR RTL  (botao PRISM, em paralelo)
     verilogSyntaxCheck -> [iverilog -tnull] ->
     generateProjectHierarchy -> [yosys] read_verilog+hierarchy+proc+
        write_json -> Temp/project_hierarchy.json -> arvore de modulos
     PRISM -> [yosys] netlist JSON -> netlistsvg (in-process) -> SVG
=========================================================================
\end{lstlisting}

\subsection{Artefatos e responsáveis por estágio}

\begin{itemize}[leftmargin=1.4em]
  \item \textbf{(1)} \file{.spf} (canônico, via \code{SpfStore.read} em
        \code{loadConfig}); pastas por \code{ensureDirectories}.
  \item \textbf{(2)} \file{<proc>.cmm}/\file{.cpp}, top-level, testbench;
        \code{ensureChegueiToaqui} só no ramo Verilator.
  \item \textbf{(3a) cmmcomp} (\code{buildCmmSpec}): \file{.asm},
        \file{cmm\_log.txt}, \file{pc\_<proc>\_mem.txt}, \file{.mif}.
        \textbf{(3b) appcomp} (\code{buildAsmPreSpec}). \textbf{(3c) asmcomp}
        (\code{buildAsmSpec}): \file{Hardware/<proc>.v}, \file{.mif},
        \file{<proc>\_tb.v}.
  \item \textbf{(4)} \code{getSimulator()} (Icarus default; Verilator opt-in);
        cocotb se testbench é \file{.py}. Orquestra \code{runGtkWave}.
  \item \textbf{(5)} \code{instrumentTestbench} grava \file{instr\_<tb>.v}
        idempotente; original nunca tocado.
  \item \textbf{(6)} \code{getViewer()}: writers \file{gtkw\_proc\_writer.js} /
        \file{surfer\_layout\_writer.js}.
  \item \textbf{(7)} \file{prism.js}; \tool{netlistsvg} roda in-process (fora da
        allowlist).
\end{itemize}

\textbf{Allowlist} (basename + diretório legal): \file{cmmcomp.exe},
\file{appcomp.exe}, \file{asmcomp.exe}, \file{comp2gtkw.exe}
(\path{components/bin}); \file{iverilog.exe}, \file{vvp.exe},
\file{verilator.exe}, \file{yosys.exe}, \file{g++.exe}, \file{make.exe},
\file{perl.exe} (\path{Packages/msys/mingw64/bin}); \file{gtkwave.exe},
\file{fst2vcd.exe} (fork \tool{gtkwave-nipscern}); \file{surfer.exe}.
